\chapter*{Introduction générale}
\addcontentsline{toc}{chapter}{Introduction générale}

Le marché de la pièce automobile d'occasion répond à un double besoin : réduire le coût de l'entretien et prolonger la durée de vie de composants encore utilisables. Pourtant, la recherche reste difficile. Une même pièce peut porter plusieurs références, convenir à plusieurs véhicules et dépendre de critères précis tels que l'année, le moteur, le carburant ou la position sur le véhicule. Une annonce incomplète ou une interprétation approximative peut donc mener à un achat inutile.


Le projet E-biessa a été conçu pour structurer ce marché et simplifier la mise en relation locale. Il ne s'agit pas seulement d'un catalogue : la plateforme associe des fonctions de place de marché, de confiance, de communication et d'intelligence artificielle. L'acheteur décrit son besoin avec ses propres mots, le système extrait les contraintes, recherche les annonces sémantiquement proches, puis vérifie la compatibilité déclarée avant d'afficher une réponse.


La problématique générale est la suivante : comment construire une place de marché automobile intuitive et riche, tout en garantissant que l'intelligence artificielle ne présente jamais une compatibilité non vérifiée ? La réponse retenue repose sur une séparation nette des responsabilités. Les données transactionnelles et l'autorisation restent dans Spring Boot et MySQL. La recherche intelligente évolue dans un service Python isolé. Les règles déterministes ont toujours priorité sur la sortie du modèle.


Le premier chapitre présente le contexte et la démarche. Le deuxième formalise les besoins et les cas d'utilisation. Le troisième détaille l'architecture, les classes et les principaux échanges. Le quatrième décrit la réalisation, y compris les aides VIN, prix, image et voix. Le cinquième distingue les résultats des tests automatisés, l'évaluation du modèle et les limites expérimentales.
